\documentclass[12pt]{article}
\usepackage[utf8]{inputenc}
\usepackage[T1]{fontenc}
\usepackage{hyperref}
\title{Fayu Language Resources for K-12 Curriculum Development}
\author{Compiled by Claude Sonnet 5.5 (claude-sonnet-5-5) for Andr\'as Kornai}
\date{}
\begin{document}
\maketitle
\section{Introduction}
Fayu (Sehudate, ISO 639-3 fau) is a Lakes Plain language of Papua, Indonesia, with about 1,400 speakers. Documentation found is academic only.
\section{Primary Educational Resources}
No primers, textbooks or school materials were found.
\section{Linguistic and Reference Materials}
Page Maitland (2020), \emph{A grammatical description of Fayu: a Western Lakes Plain language of West Papua}, University of Newcastle MA thesis: \url{https://glottolog.org/resource/reference/id/580998}. Duane Clouse (1997) comparative Lakes Plain wordlist including Fayu: \url{https://transnewguinea.org/source/clouse-1997?page=47}. Overview: \url{https://en.wikipedia.org/wiki/Fayu_language}.
\section{Religious and Community Resources}
None found.
\section{Digital and Audio Resources}
None found.
\section{Curriculum Development Context}
The language is not known to be taught in schools, and no dictionary was found. Curriculum work would depend on the Maitland grammar and Clouse wordlist and on new community-led documentation.
\end{document}